\begin{abstract}
The growing imbalance between the volume of court cases and the capacity of the judicial system has increased interest in automated methods that can assist legal analysis and outcome prediction. In the Indian setting, this problem is especially challenging because judgments are long, noisy, often multilingual, and frequently spread across multiple hearings.  Prediction pipelines using text embeddings risk leaking the final decision through operative text and identity shortcuts. This paper presents a graph embedding based framework for Legal Judgment Prediction over Indian court decisions. The approach constructs a heterogeneous knowledge graph from raw judgments using legal named entities, rhetorical-roles of texts and outcome labels after multi-hearing case consolidation. The approach caters to multilingual documents. By representing cases alongside parties, courts, statutes, provisions, precedents, and argument structures, the model captures both the internal logic of a case and the shared legal relationships across the corpus. Additionally, we introduce a comprehensively annotated dataset of over 73,000 Indian court cases spanning five legal domains, which will be released to support downstream legal NLP research. The proposed framework, using a Heterogeneous Graph Transformer, achieves 80\% accuracy and a 0.795 macro-F1 score in the pooled cross-domain setting. The results demonstrate that graph-based representations outperform text-only baselines, offering a robust, interpretable alternative that successfully resists identity-based leakage shortcuts.
\end{abstract}
